\section{Exact state and transition counts}

A singleton class has two admissible statuses and a repeated class has three.  Independence of these status choices gives exactly
\[
 |V|=2^s3^r.
\]

For a fixed singleton class, its atomic transition is available in exactly half of the singleton-status assignments and in every repeated-class assignment, contributing $2^{s-1}3^r$ edges.  Summing over the $s$ singleton classes gives
\[
 s2^{s-1}3^r.
\]

For a fixed repeated class there are two event types, start and finish.  For either event, the class must be in one specified one of its three statuses, while every other repeated class has three choices and every singleton has two choices.  Thus each event type contributes $2^s3^{r-1}$ edges for that class.  Summing start and finish events over all $r$ repeated classes gives
\[
 2r2^s3^{r-1}.
\]
Therefore
\[
 |E|=s2^{s-1}3^r+2r2^s3^{r-1}.
\]
Each state has at most one currently legal next event per class, hence out-degree at most $k$ and
\[
 |E|\le k|V|\le k3^k.
\]

After all $2^k$ values of $\rho$ are available, the rank identity requires only table accesses and integer arithmetic for a state.  The bottleneck relaxation is constant work per edge apart from reconstruction bookkeeping.  Consequently the combinatorial DP phase is $O(k3^k)$; the cost of constructing the subset-rank table is a separate preprocessing term determined by the chosen finite-field linear-algebra representation.
